\section*{Chapter \ref{ch:b3:hydrogeology} --- Hydrogeology}

\begin{solution}{exo:b3:hydrogeology:1}
First screen: $z = 120 - 40 = \qty{80}{m}$, $h = 120 - 12.5 =
\qty{107.5}{m}$, pressure head $107.5 - 80 = \qty{27.5}{m}$. Second
screen: $z = \qty{40}{m}$, $h = \qty{105.0}{m}$, pressure head
\qty{65.0}{m}. The head is higher at the shallow screen, so the water
flows downward, under a gradient $(107.5 - 105.0)/40 = \num{0.0625}$:
the piezometers stand in a recharge area. The pressure alone would have
suggested the opposite.
\end{solution}

\begin{solution}{exo:b3:hydrogeology:2}
$q = K\Delta h/L = 12\times0.30/0.50 = \qty{7.2}{m/d}$. The area is
$\pi\times0.05^2 = \qty{7.85e-3}{m^2}$, so $Q = \qty{0.0565}{m^3/d}$, or
$56.5/24 = \qty{2.36}{L/h}$. The \omterm{def:b3:hydrogeology:seepage}{seepage velocity} is $7.2/0.35 =
\qty{20.6}{m/d}$.
\end{solution}

\begin{solution}{exo:b3:hydrogeology:3}
$K = 2.5/86400 = \qty{2.894e-5}{m/s}$, and $k = K\mu/(\rho g) =
\num{2.894e-5}\times\num{1.306e-3}/(999.7\times9.81) =
\qty{3.85e-12}{m^2}$. At \qty{20}{\celsius}, $K$ is multiplied by
$(998.2/\num{1.002e-3})/(999.7/\num{1.306e-3}) = \num{1.30}$: $K =
\qty{3.25}{m/d}$. The change is almost entirely the viscosity's.
\end{solution}

\begin{solution}{exo:b3:hydrogeology:4}
$T = 2.5\times180 = \qty{450}{m^2/d}$. The expansion of the water gives
$\num{9e-7}\times180 = \num{1.62e-4}$, about \qty{80}{\percent} of $S =
\num{2e-4}$; the compression of the limestone supplies the rest.
\end{solution}

\begin{solution}{exo:b3:hydrogeology:5}
By \cref{prop:b3:hydrogeology:thiem}, $T = Q\ln(r_2/r_1)/(2\pi(s_1 -
s_2)) = 1500\times\ln10/(2\pi\times1.53) = \qty{359}{m^2/d}$, and $K =
359/40 = \qty{9.0}{m/d}$.
\end{solution}

\begin{solution}{exo:b3:hydrogeology:6}
$Q/4\pi T = 1200/(4\pi\times300) = \qty{0.318}{m}$. After 1 day, $u =
300^2\times\num{5e-4}/(4\times300\times1) = \num{0.0375}$ and $W = -0.5772
+ 3.2834 + 0.0375 - 0.0004 = \num{2.743}$: $s = \qty{0.873}{m}$. After 10
days, $u = \num{0.00375}$, $W = -0.5772 + 5.5860 + 0.0038 = \num{5.013}$:
$s = \qty{1.596}{m}$. Ten times longer gives less than twice the
drawdown.
\end{solution}

\begin{solution}{exo:b3:hydrogeology:7}
$u \le 0.05$ when $t \ge r^2S/(4T\times0.05)$. At \qty{60}{m}: $3600\times
\num{2e-4}/90 = \qty{0.008}{d}$, \qty{11.5}{min}. At \qty{500}{m}:
\qty{0.556}{d}, \qty{800}{min}, over 13 hours. The time grows as $r^2$:
in a test of a day or two, the observation wells used for the straight
line must be within a few hundred metres of the pumped well.
\end{solution}

\begin{solution}{exo:b3:hydrogeology:8}
Midway, each well is \qty{200}{m} away: $u = 200^2\times\num{2e-4}/(4
\times450\times10) = \num{4.44e-4}$, $W = \num{7.142}$, $s = 0.382\times
7.142 = \qty{2.728}{m}$ from each, \qty{5.456}{m} in all. At
\qty{400}{m}, $u = \num{1.78e-3}$, $W = \num{5.757}$, $s =
\qty{2.199}{m}$. Each well lowers the water in the other by
\qty{2.2}{m}: its pump must lift the water \qty{2.2}{m} higher, which
costs energy, and the drawdown left before the pump intake runs dry is
smaller.
\end{solution}

\begin{solution}{exo:b3:hydrogeology:9}
Four streamlines between the dam base and the floor bound $N_f = 5$
channels; nine equipotentials divide the \qty{10}{m} into $N_d = 10$
drops. $q = KHN_f/N_d = 5\times10\times5/10 = \qty{25}{m^3/d}$ per metre
of dam, and \qty{7500}{m^3/d} under \qty{300}{m}. The computed shape
factor, 0.512 instead of 0.5, would give \qty{25.6}{m^3/d} per metre.
\end{solution}

\begin{solution}{exo:b3:hydrogeology:10}
$k = na^2/8 = 0.30\times(\num{5e-5})^2/8 = \qty{9.4e-11}{m^2}$, and $K =
k\rho g/\mu = \num{9.375e-11}\times998.2\times9.81/\num{1.002e-3} =
\qty{9.2e-4}{m/s}$, \qty{79}{m/d}, inside the range of clean sand, 0.3 to
\qty{150}{m/d}. The model leaves out the tortuosity of the paths, longer
than the straight tubes (it lowers $K$), the throats between the pores,
narrower than the pores (lower), the pores that end blind and carry no
flow (lower), and the spread of the pore sizes: since a tube carries a
flow in proportion to $a^4$, the widest paths carry most of the water and
a spread raises $K$ above that of tubes of the mean size.
\end{solution}

\begin{solution}{exo:b3:hydrogeology:11}
(a)~The logarithm vanishes when $2.25\,Tt = r^2S$, i.e. at $r_0 =
\sqrt{2.25\,Tt/S}$. (b)~Confined and unconfined:
\[
r_0 = \sqrt{\frac{2.25\times450\times30}{\num{2e-4}}} = \qty{12.3}{km},
\qquad r_0 = \sqrt{\frac{2.25\times450\times30}{0.2}} = \qty{390}{m},
\]
$\sqrt{1000} = 31.6$ times smaller. (c)~The pressure
spreads with the diffusivity $T/S$, enormous when $S$ is small: within
the test the cone reaches faults, the outcrops of the aquifer and rivers,
and the drawdown leaves the straight line, which tells the
hydrogeologist that a boundary has been met.
\end{solution}

\begin{solution}{exo:b3:hydrogeology:12}
(a)~Seen from the frame moving with the water, the front is spread with
standard deviation $\sigma_x = \sqrt{2Dt} = \sqrt{2\alpha_Lvt}$. It passes
$x$ around $t = x/v$, when $\sigma_x = \sqrt{2\alpha_Lx}$; at the speed
$v$, a spread of $\sigma_x$ in distance is one of $\sigma_t = \sigma_x/v$
in time. (b)~The \qty{16}{\percent} and \qty{84}{\percent} times lie one
standard deviation either side of the mean: $\sigma_t = (247.2 -
152.6)/2 = \num{47.3}\text{ days}$. Then $\alpha_L = v^2\sigma_t^2/(2x) =
0.25\times47.3^2/200 = \qty{2.8}{m}$. (The curve was computed with
\qty{3}{m}: the normal approximation slightly underestimates it.) (c)~$R =
1 + 1.9\times0.32/0.30 = \num{3.03}$: the solute travels at a third of the
water's speed.
\end{solution}

\begin{solution}{pb:b3:hydrogeology:1}
\textbf{1.} Shale $363.2 - 282.5 = \qty{80.7}{m}$; limestone $282.5 -
102.5 = \qty{180}{m}$.

\textbf{2.} $h = 363.2 - 45.0 = \qty{318.2}{m}$; pressure head at the top
$318.2 - 282.5 = \qty{35.7}{m}$. The water rises \qty{35.7}{m} above the
top of the limestone, held down by the shale: the aquifer is confined.

\textbf{3.} From $h_2 - h_1 = 550a - 100b = -1.62$ and $h_3 - h_1 = 250a
- 600b = -1.44$: $a = \num{-2.715e-3}$ and $b = \num{1.269e-3}$ per
metre.

\textbf{4.} $|\nabla h| = \sqrt{a^2 + b^2} = \num{0.0030}$. The water
flows along $-\nabla h = (\num{2.715e-3}, \num{-1.269e-3})$, east and
south: azimuth $\ang{180} - \arctan(2.715/1.269) = \ang{115}$.

\textbf{5.} The points before about \qty{10}{min} lie above the line: there
$u > 0.05$ and the Theis curve is above its logarithmic approximation
($W > -\gamma - \ln u$, as the figure of the well function shows).

\textbf{6.} The best line through the points from \qty{15}{min} on passes
\qty{1.75}{m} at \qty{100}{min} and \qty{2.63}{m} at \qty{1000}{min}:
$\Delta s = \qty{0.87}{m}$ per log cycle.

\textbf{7.} $T = 2.3\times2160/(4\pi\times0.872) = \qty{453}{m^2/d}$.

\textbf{8.} The line falls \qty{0.87}{m} per cycle from \qty{1.75}{m} at
\qty{100}{min}: zero is reached $1.75/0.872 = 2.01$ cycles earlier,
$t_0 = \qty{0.98}{min} = \qty{6.83e-4}{d}$.

\textbf{9.} $S = 2.25\,Tt_0/r^2 = 2.25\times453\times\num{6.83e-4}/3600 =
\num{1.9e-4}$, within the range of \omterm{def:b3:hydrogeology:aquifer-types}{confined aquifers}.

\textbf{10.} At \qty{15}{min}, $u = 3600\times\num{1.9e-4}/(4\times453
\times0.0104) = \num{0.036}$, and less afterwards.

\textbf{11.} $K = 453/180 = \qty{2.5}{m/d}$, a fractured limestone.

\textbf{12.} $u = 3600\times\num{2e-4}/(4\times450\times2) =
\num{2.0e-4}$, $W = -0.5772 + 8.5172 + 0.0002 = \num{7.940}$, $s =
0.382\times7.940 = \qty{3.03}{m}$, as recorded.

\textbf{13.} $u = 10^6\times\num{2e-4}/(4\times450\times30) =
\num{0.0037}$, $W = \num{5.02}$, $s = \qty{1.92}{m}$.

\textbf{14.} $r_0 = \sqrt{2.25\times450\times30/(\num{2e-4})} =
\qty{12.3}{km}$, more than the width of the region. The cone has reached
the edges of the limestone, its outcrops and the fault, which may set it
against less permeable rock: the infinite aquifer of Theis no longer
describes it, and only the record of the drawdowns will say whether a
boundary feeds the well or stops the flow.

\textbf{15.} $2160\times30 = \qty{64800}{m^3}$. It comes from storage:
the water released by the expansion of the water and the compression of
the limestone, over the whole area of the cone. Integrating $S\,s$ over
the cone gives back $Qt$.

\textbf{16.} $u = 10^6\times0.18/(4\times450\times30) = \num{3.33}$, $W
= \num{0.0086}$, $s = \qty{3}{mm}$: the farm would not notice. The
pumping would draw on the pores near the well instead of the pressure
of the whole aquifer.

\textbf{17.} $q = 2.5\times0.0030 = \qty{0.0075}{m/d}$; $v = 0.0075/0.18
= \qty{0.042}{m/d}$, \qty{15}{m} a year.

\textbf{18.} $200/0.0417 = \num{4800}\text{ days}$, \num{13.1}~years; the
pollutant, \num{32.9}~years.

\textbf{19.} Through a cylinder of radius $r$ the flow $Q$ crosses the
area $2\pi rb$ of which the pores are the fraction $n_e$: $\mathrm dr/\mathrm
dt = -Q/(2\pi rbn_e)$, and integrating $2\pi bn_er\,\mathrm dr = -Q\,
\mathrm dt$ from $r$ to 0 gives $t = \pi n_ebr^2/Q$. Here $\pi\times0.18
\times180\times200^2/2160 = \num{1885}\text{ days}$, \num{5.2}\text{ years} for the
water and $2.5\times5.2 = \num{12.9}\text{ years}$ for the pollutant: the well
brings it more than twice as soon.

\textbf{20.} $\alpha_L = \qty{20}{m}$ and $\sqrt{2\times20\times200} =
\qty{89}{m}$. The concentration reaches \qty{16}{\percent} of its
final value \qty{89}{m} ahead of the mean front: that level reaches the
well when the mean front is still \qty{89}{m} out, at $(111/200)^2\times12.9 = \num{4.0}\text{ years}$, a third
of the mean time. The first traces matter for drinking water, not the
mean.

\textbf{21.} For: the well captures the pollutant, so it cannot spread
further down-gradient, and the water can be treated where it is pumped.
Against: the pumping draws the pollutant into a supply well sooner than
it would otherwise travel anywhere, and the water must be monitored and
treated for decades.

\textbf{22.} The \omterm{def:b3:hydrogeology:storage}{transmissivity} of the Varden limestone is $T \approx
\textbf{\qty[text-series-to-math=true]{450}{m^2/d}}$ (\qty{453}{m^2/d}
from the test).
\end{solution}
